\section{Uniform recovery and exact bands at wild indices}\label{sec:all-wild-indices}
This section combines the pointwise graph bound with an equivariant cyclic-point construction whose number of degree-$p$ stages is prescribed by the divisibility of the base index.

\subsection{A graph bound for the full filtration class}
The determinant first fixes the total dimension; the mixed and supersingular bounds then exhaust the remaining comparisons.

\begin{proposition}\label{prop:graph-uniform-full}
Let $p\ge5$ and $n\ge2$. In the standing height-two filtration class, $e_K\le(p-1)p^{n-2}$ implies recovery of the special character from actual $p^n$-torsion. Consequently
\[
 \nuMin(p,e,2)\le1+\left\lceil\log_p\frac{e}{p-1}\right\rceil.
\]
This is a consequence of Theorem~\ref{thm:all-index-upper}; the following direct proof uses the pointwise graph valuation bound.
\end{proposition}
\begin{proof}
Pass to the completed maximal unramified base. This preserves $e=e_K$, inertia and the actual torsion comparison. For an allowed height-two group the determinant on inertia is $\chi_p^d\epsilon$, where $0\le d\le2$ is the total dimension of its potentially good graded pieces and $\epsilon$ is finite of order dividing $p-1$, as in the proof of Theorem~\ref{thm:first-wild-general}.

The cyclotomic image modulo $p^n$ is nontrivial. Indeed a primitive $p^n$th root satisfies
\[
 v_K(\zeta_{p^n}-1)=\frac{e}{(p-1)p^{n-1}}\le\frac1p<1,
\]
where the value follows from the Eisenstein polynomial $\Phi_{p^n}(1+X)$ over a Cohen coefficient ring. It cannot belong to the integer-valued field after unramified passage. Since $\mu_p\subset K$, that nontrivial image is a $p$-group, and restriction to wild inertia has the same image because the tame quotient has no nontrivial finite $p$-quotient. The finite determinant factors become trivial on wild inertia. Equality of determinants modulo $p^n$ therefore implies $p\mid d-d'$, and $|d-d'|\le2<p$ gives $d=d'$. Reduction modulo $p$ then identifies $\epsilon$ and $\epsilon'$ by injectivity on $\mu_{p-1}$.

If $d=0$, the special sum character is the Brauer lift of the residual semisimplification; the common $p$-torsion determines it. Cartier duality gives $d=2$. If $d=1$, each group is either whole potentially supersingular or has a two-step potentially good height-one filtration. In the ordinary case the unique connected--\'etale filtration descends, giving the second possibility. Two height-one-filtered groups compare residually: their special characters are sums of the finite factors, and $\chi_p\equiv1\pmod p$.

Put $v=v_p(e)$. The assumed bound gives
\[
                    e\le(p-1)p^{n-2}<p^{n-1},\qquad v\le n-2.
\]
A mixed comparison would require, by Proposition~\ref{prop:mixed-obstruction},
\[
                        e>p^{2n-3-v}\ge p^{n-1},
\]
which is impossible. Two whole supersingular groups with unequal special characters would require, by Proposition~\ref{fc-thm:valuation},
\[
                         e(p^n-1)\ge(p-1)p^{2n-2}.
\]
But
\[
 \begin{aligned}
 e(p^n-1)&\le(p-1)p^{n-2}(p^n-1)\\
         &=(p-1)\bigl(p^{2n-2}-p^{n-2}\bigr)
          <(p-1)p^{2n-2}.
 \end{aligned}
\]
This excludes the last case, including the endpoint of the asserted range.

For the logarithmic formulation let $t$ be the least nonnegative integer with $e/(p-1)\le p^t$. Applying the result with $n=t+2$ gives $\nuMin\le t+1$. This is an integer-threshold statement; no approximation to a logarithm is needed.
\end{proof}

\subsection{Prescribing the degree-$p$ stages}
The next construction produces actual groups at each index, including those with an arbitrarily large $p$-part.

\begin{proposition}\label{prop:mixed-stage-construction}
Let $p\ge5$, $n\ge2$ and $0\le t\le n-2$. Suppose $e$ is a positive multiple of $(p-1)p^t$ and
\[
                             e>p^{2n-3-t}.
\]
There are actual groups $M,H$ over a complete field of residue $\overline{\F}_p$, exact absolute index $e$ and containing $\mu_p$, with $M[p^n]\simeq H[p^n]$ and different special characters. The group $M$ is whole potentially supersingular, while $H$ has a two-step potentially good height-one filtration.
\end{proposition}
\begin{proof}
Set
\[
 q=p^2,\qquad N=q-1,\qquad j=n-1-t\ge1,\qquad d=Ne,
 \qquad x_0=\frac e{p^{2n-2-t}}.
\]
Choose an integer $A$ with $x_0<A<px_0$ and $N\nmid A$. Such a choice exists because $px_0>1$. If $x_0<1$, take $A=1$; otherwise take $\lfloor x_0\rfloor+1$ and increase it by one if divisible by $N$. The result is at most $x_0+2<px_0$, since $p\ge5$.

Put
\[
                       b=\frac{Ne}{p^t}-(p-1)Ap^{2j}.
\]
It is integral, and the two strict comparisons are
\begin{align}
 b-(q-p)Aq^{j-1}
   &=\frac N{p^t}\bigl(e-Ap^{2n-3-t}\bigr)>0,
       \label{eq:stage-first}\\
 (q-p)Aq^j-b
   &=\frac N{p^t}\bigl(Ap^{2n-2-t}-e\bigr)>0.
       \label{eq:stage-second}
\end{align}
In particular $b>0$. Also
\begin{equation}\label{eq:stage-descent-congruence}
                         b\equiv-(p-1)A\pmod N,
\end{equation}
because $q^j\equiv1\pmod N$ and $p^t\mid e$.

Over $W=W(\overline{\F}_p)$ use the integral height-two formal group $F_u$ and the full homogeneous identities of~\eqref{eq:homogeneous-logarithm}. Define
\[
 Y_0=X^A,\qquad Y_{i+1}=[p]_{X^b}(Y_i),\qquad
 P(X)=\Phi_{X^b}(Y_{n-1}),\qquad \Phi_u(T)=[p]_u(T)/T.
\]
These are integral power series. In the first $j$ compositions the selected degree is
\[
                             h_i=Aq^i\quad(0\le i\le j).
\]
Indeed~\eqref{eq:stage-first} says $b>(q-p)h_{j-1}$, hence the same inequality holds for all earlier inputs. It makes $T^q$ strictly earlier than $uT^p$. Intermediate coefficients contain $u$ modulo $p$, so are later; degrees above $q$ are later by integrality.

For the remaining $t$ compositions, put $h_{i+1}=b+ph_i$, starting from $h_j=Aq^j$. Inequality~\eqref{eq:stage-second} gives $b<(q-p)h_j$, and the $h_i$ strictly increase. Thus $uT^p$ is strictly first at every remaining step. All selected coefficients are one modulo $p$, so
\[
 Y_i\equiv X^{h_i}+O(X^{h_i+1})\pmod p,
 \qquad H:=h_{n-1}=p^tAq^j+b\frac{p^t-1}{p-1}.
\]
For $t=0$ the last formula is $H=Aq^j$ and there are no degree-$p$ compositions before the final divided series.

In $P$ the term $uY_{n-1}^{p-1}$ has degree
\begin{equation}\label{eq:stage-eisenstein-degree}
 b+(p-1)H=p^tb+(p-1)p^tAq^j=Ne=d.
\end{equation}
Its competitor $Y_{n-1}^{q-1}$ is later because $b<(q-p)H$, by~\eqref{eq:stage-second} and $H\ge Aq^j$. The intermediate terms have a factor $u$ modulo $p$ and are later than the selected term. Consequently
\[
                    P(0)=p,\qquad P(X)\equiv X^d+O(X^{d+1})\pmod p.
\]
Weierstrass preparation gives an Eisenstein ring $R=W[[X]]/(P)$ with uniformizer $x$ and $v_R(p)=d$. Equation~\eqref{eq:stage-eisenstein-degree} implies $H<d$, and all earlier $h_i$ are smaller. Coefficients divisible by $p$ therefore do not affect the values $v_R(Y_i(x))=h_i$. In particular $Y_{n-1}(x)\ne0$, while $\Phi(Y_{n-1}(x))=0$ makes its $p$-multiple zero. Thus $x^A$ has exact order $p^n$.

Choose $\zeta\in W$ of order $N$ and set $a=\zeta^{-A}$. Congruence~\eqref{eq:stage-descent-congruence} and the exact homogeneous identities show that $X\mapsto\zeta X$ carries $u$ to $a^{p-1}u$, every $Y_i$ to $a^{-1}Y_i$, and $P$ to itself. Hence $\sigma(x)=\zeta x$ defines an automorphism of $R$ of exact order $N$. For $L=\operatorname{Frac}R$ and $K=L^{\langle\sigma\rangle}$, the extension $L/K$ is totally tame of degree $N$. Thus $e_K=d/N=e$ and the residue field is $\overline{\F}_p$.

The full formal-group isomorphism $\sigma^*F\to F$, $T\mapsto aT$, satisfies the cocycle identity and fixes the marked point because $a\sigma(x^A)=x^A$. Descent of the compatible finite generic levels gives an actual height-two group $G/K$ with that rational point. Its special eigenvalues are $a,a^p$, and $a\ne1$ because $N\nmid A$. This suffices; $a$ need not have exact order $N$.

The field $K$ contains $\mu_p$: write $-p=\pi_K^eU$, extract a $(p-1)$st root of the unit $U$ by Hensel's lemma, and use $(p-1)\mid e$. The resulting root of $-p$ generates the unique tame degree-$(p-1)$ extension of the Cohen field, namely its cyclotomic extension.

Finally Lemmas~\ref{lem:point-ambiguity} and~\ref{lem:extension-lift} supply actual $M,H$. Concretely, quotient $G$ by the marked constant cyclic subgroup and twist by the inverse finite determinant factor $\psi^{-1}$. The finite extension is then an extension of $(\Z/p^n)(\psi^{-1})$ by $\mu_{p^n}$. Compatible Kummer lifts and prime-to-$p$ averaging give the height-one filtered group $H$ with that same truncation. Its special eigenvalues are $1,a^{-(p+1)}$, while those of $M$ are $a^{-1},a^{-p}$. Their trace difference
\[
                         -(a^{-1}-1)(a^{-p}-1)
\]
is nonzero. This proves the exact field and actual-group assertions.
\end{proof}

\subsection{The sharp mixed criterion and full-class bands}
The number of degree-$p$ stages can be chosen from $v_p(e)$, and the uniform recovery range makes the resulting lower examples sharp on explicit infinite bands.

\begin{proof}[Proof of Theorem~\ref{thm:mixed-all-indices}]
Necessity is Proposition~\ref{prop:mixed-obstruction}. For sufficiency put $v=v_p(e)$ and $t=\min(v,n-2)$. If $v\le n-2$, the desired strict inequality is exactly the hypothesis of Proposition~\ref{prop:mixed-stage-construction}. If $v\ge n-1$, then
\[
                 e\ge(p-1)p^v\ge(p-1)p^{n-1}>p^{n-1}
                    =p^{2n-3-t},
\]
so that proposition again applies. In both cases $p^t(p-1)\mid e$, and its construction has all the required properties.
\end{proof}

\begin{proof}[Theorem~\ref{thm:wild-bands-exact} for $2\le k\le p+1$]
In this paragraph assume $2\le k\le p+1$. Put $e=kp^r(p-1)$ and $n=r+2$. Use Proposition~\ref{prop:mixed-stage-construction} with $t=r=n-2$. Its strict threshold is $e>p^{r+1}$, which follows from $k(p-1)>p$. It supplies an actual mixed pair with isomorphic $p^{r+2}$-torsion and unequal special characters, so $\nuMin\ge r+2$.

The construction permits $A=1$ uniformly, since
\[
                \frac{k(p-1)}{p^2}<1<\frac{k(p-1)}p.
\]
Its deformation exponent is $b=(p-1)((k-1)p^2-k)>0$, independent of $r$. There is one degree-$p^2$ stage, then $r$ degree-$p$ stages, before the final divided series. Thus each index is attained directly, without requiring that a lower example survive a prescribed base change.

For $k\le p$, Proposition~\ref{prop:graph-uniform-full} at level $r+3$ gives the upper bound, since $e\le(p-1)p^{r+1}$. For the entire range $k\le p+1$, Theorem~\ref{thm:all-index-upper} applies because
\[
 e\le(p^2-1)p^r<p^{r+2}.
\]
Thus $\nuMin\le r+2$ in this range, proving equality. In particular this endpoint is $\nuMin(p,(p^2-1)p^r,2)=r+2$. The corresponding finite band follows by $r=1$, and the diagonal formula follows by $k=p$, $r=s-1$. At $k=p$, the valuation of $e$ is $r+1$ rather than $r$; the direct construction above remains valid because it requires only $(p-1)p^r\mid e$. The larger coefficient interval, with its distinction between coefficients prime to $p$ and divisible by $p$, is treated in Section~\ref{sec:extended-wild-bands}.
\end{proof}
